\documentclass[10pt,a4paper]{amsart}
\usepackage[margin=19mm]{geometry}
\usepackage{amsmath,amssymb,mathtools}
\usepackage[scr=rsfs,cal=euler]{mathalfa}
\usepackage{xfrac}
\usepackage[unicode,hidelinks,bookmarks=false]{hyperref}
\newcommand{\ie}{i.e.\ }
\newcommand{\cf}{cf.\ }
\newcommand{\resp}{respectively\ }
\newcommand{\Subsection}{\subsection}
\providecommand{\nobreakdash}{\nobreak\hbox{-}}
\setlength{\emergencystretch}{2em}
\multlinegap0pt
\usepackage{bm}
\usepackage{bbm}
\usepackage{dsfont}
\usepackage{xspace}
\usepackage{cancel}
\usepackage{mathtools}
\usepackage{amsthm}
\makeatletter
\@ifundefined{thm}{\newtheorem{thm}{Theorem}}{}
\@ifundefined{prop}{\newtheorem{prop}[thm]{Proposition}}{}
\makeatother
\title[Optimal Dual Frame Pairs: A Synergy with Graph Theory]{Optimal Dual Frame Pairs: A Synergy with Graph Theory}
\author{Shankhadeep Mondal, Ram Narayan Mohapatra}
\date{Source: arXiv:2507.23249v1 (CC BY 4.0)}
\begin{document}
\maketitle

\noindent\emph{Source and licence.} Optimal Dual Frame Pairs: A Synergy with Graph Theory. Version arXiv:2507.23249v1 (CC BY 4.0), distributed under the Creative Commons Attribution 4.0 International licence, as recorded in the arXiv licence metadata for this exact version and shown on the abstract page https://arxiv.org/abs/2507.23249v1. Full rights notice: licenses/arxiv-2507.23249-CC-BY-4.0.txt.\par\smallskip

\noindent\textit{Editorial scope.} Counted unit: the Thm whose source label is \texttt{None} in the article (source0.tex lines 461--463, source label \texttt{None}), delivered together with the article's own complete original proof of it (source0.tex lines 464--490, one \texttt{proof} environment of 27 lines). No mathematics is written, repaired or re-derived: statement and proof are the article's own text line for line. Required context retained uncounted: equation2 (lines 168--170); prop3point1 (lines 173--175); equation6 (lines 323--325); prop4point1 (lines 328--330); thm5point2 (lines 410--412). Every definition, display and label of those lines is retained unchanged. Omitted from this package and not redistributed: 1--167, 171--172, 176--322, 326--327, 331--409, 413--460, 491--701. Nothing inside a retained excerpt is omitted or altered: every retained body is delivered exactly as the article's lines are written, so no omission receipt of any kind accompanies this package. Citations: the retained excerpts carry no citation command at all, so no bibliography entry of the article is transcribed and no reference is redistributed. Reference adaptation: none was needed and none was made; every reference inside the delivered unit resolves inside it, so no unresolved reference or citation is produced. Printed slips kept literal: no printed word, symbol, hypothesis or number of the counted statement or of its proof was altered. Layout and class: the article's own class is \texttt{elsarticle} with packages including amssymb, amsmath, amsthm, inputenc, marginnote, amsfonts, textcomp, color, soul, placeins, tikz, relsize, makecell, graphicx; the wrapper is the lane's pinned \texttt{amsart} editorial wrapper at 10pt on A4, which restates the theorem-like environments the retained text needs and adds the article's own macro definitions that the retained text needs (none), with extra wrapper packages (bm, bbm, dsfont, xspace, cancel, mathtools). The delivered numbering is the unit's own. Assets: no retained span contains a figure environment, a graphics-inclusion command, a PDF-page inclusion, a table or a TikZ picture; no image file is copied into this package, the package declares no asset dependency and no image file is redistributed with it. Licence: version arXiv:2507.23249v1 of this article is distributed under the Creative Commons Attribution 4.0 International licence, as recorded in the arXiv licence metadata for this exact version and shown on the abstract page https://arxiv.org/abs/2507.23249v1. A full rights notice is delivered at licenses/arxiv-2507.23249-CC-BY-4.0.txt. Characters: every retained excerpt is pure ASCII, so the delivered engine, XeLaTeX with its default Latin Modern text font, reports no missing character.
	\begin{equation}\label{equation2}
		\mathcal{E}_{1}^p (F,G) = \left\{\frac{1}{N} \sum\limits_{i=1}^N \left|\langle f_i, g_i \rangle \right|^p\right\}^{\frac{1}{p}}
	\end{equation}\\

	\begin{prop}\label{prop3point1}
		The  value of $ \delta_{1}^p$ is $\frac{n}{N}$. Moreover, an $(N,n)$ dual pair $(F,G) $ is a $1-$erasure spectrally optimal  if and only if it is $1-$uniform.
	\end{prop}

\begin{equation}\label{equation6}
	\mathcal{O}_{1}^p (F,G) = \left\{\frac{1}{N} \sum \limits_{i=1}^N\|f_i\|^p\,\| g_i \| ^p\right\}^{\frac{1}{p}}.
\end{equation}\\

	\begin{prop}\label{prop4point1}
	Let $\mathcal{H}$ be an $n$ dimensional Hilbert space. Then the  value of $ \Delta_{1}^p$ is $\frac{n}{N}$ and  $ \mathcal{O}_{1}^p = \left\{(F,G): \frac{1}{N} \sum\limits_{i=1}^N \left(\|f_i\|\,\| g_i \| \right)^p =  \left(  \frac{n}{N}  \right)^p \right\}$ 
\end{prop}

\begin{thm} \label{thm5point2}
	Let $\Gamma$ be a connected graph of $N$ vertices and $F= \{f_i\}_{i=1}^N$ is a $\mathcal{L}_{\Gamma}(N,N-1)$ frame for $\mathbb{R}^{(N-1)}.$ If $F$ is a  tight frame then $\left( F, S_{F}^{-1}F \right) \in 	\mathcal{E}_{2}^p .$
\end{thm}

\begin{thm}
	Let $\Gamma$ be a graph of $N$ vertices with no null vertex and $F= \{f_i\}_{i=1}^N$ is a $\Gamma(N,n)$ frame for $\mathbb{C}^{(n)}.$ If $F$ is a  tight frame then $\left( F, S_{F}^{-1}F \right) $ is both $1-$erasure optimal dual pair under spectral radius and operator norm. 
\end{thm}

\begin{proof}
	Let $F$ be a tight frame with bound $A$ and $\Gamma_1, \Gamma_2,\ldots, \Gamma_{(N-n)}$ are the connected component of $\Gamma.$ Let the number of vertices in $\Gamma_i$ is $n_i.$ The Gramian matrix of $F$ is $$\mathcal{G} = \mathcal{L}(\Gamma) = \begin{pmatrix}
		\mathcal{L}(\Gamma_1) & 0  & \cdots &0 &0\\
		0 & \mathcal{L}(\Gamma_2)  & \cdots &0 &0 \\
		\vdots &\vdots  &\vdots &\vdots &\vdots \\
		0 & 0  & \cdots &0 &  \mathcal{L}(\Gamma_{(N-n)}) \\
		
	\end{pmatrix}.$$ 
As the frame operator of corresponding to $F$ is $S_F = AI_{n\times n},$ the eigenvalues of $\mathcal{L}(\Gamma)$ are $A$ with multiplicity $n$ and $0$ with multiplicity $(N-n).$ As each $\Gamma_i$ is connected, $\mathcal{L}(\Gamma_i)$ has rank $(n_i -1)$ and so $L(G_i) $ has eigenvalues $A$ with multiplicity $(n_i -1)$ and  $0$ with multiplicity 1. Now using the similar argument as in the proof of Theorem  \ref{thm5point2}, we can prove that each $\Gamma_i$ is regular. Let each $\Gamma_i$ is $r_i -$regular. Obviously $r_i > 0,\,$ for all $1 \leq i \leq N-n.$ Now, the adjacency matrix of $G$ is $$\mathcal{A}(\Gamma) =  \begin{pmatrix}
	\mathcal{A}(\Gamma_1) & 0  & \cdots &0 &0\\
	0 & \mathcal{A}(\Gamma_2)  & \cdots &0 &0 \\
	\vdots &\vdots  &\vdots &\vdots &\vdots \\
	0 & 0  & \cdots &0 &  \mathcal{A}(\Gamma_{(N-n)})\end{pmatrix} .$$ 
	Therefore, 
	$$ \mathcal{G} = \mathcal{L}(\Gamma) = \mathcal{D}(\Gamma) - \mathcal{A}(\Gamma) = \begin{pmatrix}
		r_{1}I - \mathcal{A}(\Gamma_1) & 0  & \cdots &0 &0\\
		0 & r_{2}I - \mathcal{A}(\Gamma_2)  & \cdots &0 &0 \\
		\vdots &\vdots  &\vdots &\vdots &\vdots \\
		0 & 0  & \cdots &0 &  r_{N-n}I - \mathcal{A}(\Gamma_{(N-n)})\end{pmatrix} .$$
	
	As $S_F = AI_{n\times n} = \Theta_{F}^* \Theta_F,$ we then have, $\mathcal{G}^2 = \Theta_F \Theta_{F}^{*} \Theta_F \Theta_{F}^{*} = \Theta_F S_F \Theta_{F}^* = A\mathcal{G}.$ This gives, $\left( r_{i} I -\mathcal{A}(\Gamma_i)    \right)^2 = A\left( r_{i} I -\mathcal{A}(\Gamma_i)    \right),\;1 \leq i \leq N-n.$ After simplifying we get 
	\begin{align}\label{eqn5point7}
		\left(\mathcal{A}(\Gamma_i)\right)^2 + (A- 2r_i)\mathcal{A}(\Gamma_i) + (r_{i}^2 - 2r_i)I) =0 , \;\;\forall i.
		\end{align} 
	 It is easy to see that each diagonal entries of each $\left(\mathcal{A}(\Gamma_i)\right)^2$ is $r_i$ and $\mathcal{A}(\Gamma_i)$ is $0$ and thus comparing the  diagonal entries on both side of \eqref{eqn5point7}, we get $r_{i} + r_{i}^2 - Ar_{i} =0,\; \forall i. $ Thus, $r_i = A -1,\; 1\leq i \leq N-n$ and hence $\Gamma$ is a $(A-1)-$regular graph. Hence, all the diagonal entries of $\mathcal{G}$ are same. This in turn gives $\|f_i\|^2$ is a constant for all $i.$ Using the fact that $n =\sum\limits_{i=1}^N \langle f_i,S_{F}^{-1}f_i \rangle = \sum\limits_{i=1}^N \frac{1}{A}\|f_i\|^2,$ we have $\|f_i\|^2 = \frac{An}{N},\; \forall i.$ Thus by \eqref{equation2}and \eqref{equation6}, we have $\mathcal{E}_{1}^p (F, S_{F}^{-1}F) = \frac{n}{N}= \mathcal{O}_{1}^p (F, S_{F}^{-1}F). $ Therefore, by Proposition\ref{prop3point1} and \ref{prop4point1} the result follows.
	
\end{proof}

\end{document}
